(1) $I$ commutes with every even element and $I^{2} = -1$ in $\Cl(1,3)$, so $(I\sigma_r)^{2} = I^{2}\sigma_r^{2} = -1$ using $\sigma_r^{2} = +1$. The eigenplane statement is \cite[Thm.~6.1]{ButlerT5}, whose spectrum is $(-2,-2,1,1,1,1)$ with the doubled eigenvalue attained exactly on $\operatorname{span}\{\sigma_r, I\sigma_r\}$. (2) With $(I\sigma_r)^{2} = -1$ the exponential series splits into the cosine and sine series, giving the displayed rotor; $R_{2\pi} = \cos\pi + I\sigma_r\sin\pi = -1$, and conjugation by $-1$ is the identity, so the conjugation action closes at $2\pi$ while the rotor itself closes at $4\pi$. (3) The two flows are the two real one-parameter subgroups of the rotor group of the plane pair, related by the substitution $\lambda = i\theta$ in the exponential; the substitution exchanges $\cosh, \sinh$ with $\cos, \sin$ exactly because the two generators differ by the factor $I$ with $I^{2} = -1$.
